\section{既約表現}

\begin{deftcolorbox}[title = 群の表現]
    群$G$の各元$g$に対して空間$V$上の可逆な線形変換を対応させる写像$D : G \rightarrow GL(V)$が存在し
    \begin{equation*}
        D(g_1)D(g_2) = D(g_1 g_2)
    \end{equation*}
    が成り立つとき$(D,V)$を\textbf{群$G$の表現}という.
\end{deftcolorbox}

\begin{deftcolorbox}[title = 不変部分空間]
    群$G$の表現$(D,V)$について, $V$の部分空間$W \in V$について
    \begin{equation*}
        \forall g \in G,\ D(g)W \in W
    \end{equation*}
    を満たすとき, $W$を\textbf{$G-$不変部分空間}という.
\end{deftcolorbox}

\begin{deftcolorbox}[title = 既約表現]
    群$G$の表現$(D,V)$について, $V$と$0$以外に$G-$不変部分空間を持たないとき
    \textbf{既約表現}という
\end{deftcolorbox}